% ══════════════════════════════════════════════════════════════════════════════
% PA-SSL: Physiology-Aware Self-Supervised Learning for ECG Anomaly Detection
% Formatted for IEEE Transactions on Medical Imaging / NeurIPS
% ══════════════════════════════════════════════════════════════════════════════

\documentclass[journal,10pt]{IEEEtran}

% ─── PACKAGES ─────────────────────────────────────────────────────────────────
\usepackage[utf8]{inputenc}
\usepackage[T1]{fontenc}
\usepackage{amsmath,amssymb,amsfonts}
\usepackage{booktabs}
\usepackage{graphicx}
\usepackage{xcolor}
\usepackage{algorithm}
\usepackage{algorithmic}
\usepackage{hyperref}
\usepackage{multirow}
\usepackage{subcaption}
\usepackage{array}
\usepackage{tabularx}
\usepackage{colortbl}
\usepackage{url}
\usepackage{cite}
\usepackage{balance}

% ─── CUSTOM COMMANDS ──────────────────────────────────────────────────────────
\newcommand{\method}{PA-SSL}
\newcommand{\ie}{i.e.,~}
\newcommand{\eg}{e.g.,~}
\newcommand{\etal}{\textit{et al.}}
\newcommand{\wrt}{w.r.t.~}
\newcommand{\figref}[1]{Fig.~\ref{#1}}
\newcommand{\tabref}[1]{Table~\ref{#1}}
\newcommand{\secref}[1]{Section~\ref{#1}}
\newcommand{\eqnref}[1]{Eq.~(\ref{#1})}

\definecolor{qrsred}{HTML}{B2182B}
\definecolor{passl}{HTML}{2166AC}

\begin{document}

% ══════════════════════════════════════════════════════════════════════════════
% TITLE & AUTHORS
% ══════════════════════════════════════════════════════════════════════════════

\title{PA-SSL: Physiology-Aware Self-Supervised Learning\\for ECG Anomaly Detection}

\author{
\IEEEauthorblockN{Venkateramanan Manivannan$^a$}
\IEEEauthorblockA{$^a$School of Computer Science and Engineering (SCOPE), Vellore Institute of Technology (VIT), Vellore, 632014, Tamil Nadu, India}
}

\maketitle

% ══════════════════════════════════════════════════════════════════════════════
% ABSTRACT
% ══════════════════════════════════════════════════════════════════════════════

\begin{abstract}
Self-supervised learning (SSL) has emerged as a powerful paradigm for learning representations from unlabeled electrocardiogram (ECG) recordings. However, existing contrastive learning approaches for ECG apply generic data augmentations (\eg time reversal, random cropping, unconstrained masking) that can \emph{destroy clinically diagnostic morphological features} such as the QRS complex, P-wave, and T-wave morphology. We propose \textbf{\method{}} (Physiology-Aware Self-Supervised Learning), a contrastive framework that introduces a library of \emph{eight physiology-constrained augmentations} jointly with \emph{multi-scale temporal adjacency contrastive learning}. Each augmentation is formally designed to (i) preserve QRS complex morphology via smooth blending masks around the R-peak, (ii) maintain physiologically plausible heart rate bounds (30--200~BPM), and (iii) inject only clinically realistic noise patterns (baseline wander, EMG artifacts, powerline interference). The temporal branch leverages the natural sequential structure of heartbeats within a recording, treating multi-scale neighboring beats ($\pm 1, \pm 2, \pm 3$) as positive pairs to encode rhythm-level information.

We evaluate \method{} on three large-scale public ECG databases — PTB-XL (21,799 recordings), MIT-BIH Arrhythmia (48 half-hour recordings), and Chapman-Shaoxing (10,646 recordings) — using patient-aware data splitting to prevent information leakage. On downstream classification and anomaly detection tasks, \method{} achieves state-of-the-art or competitive results with only 1\%--10\% labeled data, outperforming both supervised baselines and prior SSL methods. Quantitative augmentation validity metrics confirm that our physiology-aware augmentations achieve $>$0.95 QRS Pearson correlation, compared to $<$0.70 for naive augmentations. Extensive ablation studies, cross-dataset transfer evaluation, calibration analysis, and Grad-CAM interpretability demonstrate the robustness and clinical relevance of the learned representations.
\end{abstract}

\begin{IEEEkeywords}
Self-supervised learning, ECG, contrastive learning, anomaly detection, physiology-aware augmentations, temporal contrastive, cardiac signal processing.
\end{IEEEkeywords}


% ══════════════════════════════════════════════════════════════════════════════
% I. INTRODUCTION
% ══════════════════════════════════════════════════════════════════════════════

\section{Introduction}
\label{sec:introduction}

Electrocardiogram (ECG) interpretation is fundamental to cardiovascular diagnostics, with over 300 million ECGs recorded annually worldwide~\cite{Schlapfer2017}. While deep learning models have achieved expert-level performance on ECG classification~\cite{Hannun2019, Ribeiro2020}, these approaches require thousands of cardiologist-annotated recordings — a resource-intensive bottleneck that limits scalability, particularly for rare arrhythmias.

Self-supervised learning (SSL) offers a compelling alternative by learning representations from \emph{unlabeled} data. Contrastive approaches such as SimCLR~\cite{Chen2020SimCLR}, MoCo~\cite{He2020MoCo}, and BYOL~\cite{Grill2020BYOL} have demonstrated that strategic data augmentation combined with contrastive objectives can produce representations that rival or exceed supervised pre-training when subsequently fine-tuned with limited labels. Recent work has begun adapting these frameworks to ECG signals~\cite{Kiyasseh2021clocs, Mehari2023, Oh2022lead}.

However, a critical gap remains: \textbf{existing SSL augmentation strategies for ECG are physiologically naive}. Generic augmentations borrowed from computer vision (\eg random cropping, time reversal, unconstrained amplitude scaling) or time-series literature (\eg arbitrary jittering, permutation) can fundamentally corrupt the morphological features that clinicians rely upon for diagnosis:

\begin{itemize}
    \item \textbf{QRS complex destruction}: Random masking or strong scaling can obliterate the QRS complex — the primary feature for rhythm classification.
    \item \textbf{Time reversal}: Produces physiologically impossible signals where the T-wave appears before the QRS complex.
    \item \textbf{Unconstrained warping}: Can compress or stretch heartbeats to heart rates outside the 30--200~BPM physiological range.
\end{itemize}

These issues have practical consequences: SSL models trained with naive augmentations may learn \emph{spurious invariances} to non-physiological signal transformations, rather than clinically meaningful invariances to noise, electrode placement variability, and heart-rate variations.

In this paper, we introduce \textbf{\method{}}, a physiology-aware self-supervised learning framework that addresses these limitations through two key contributions:

\vspace{0.5em}
\noindent\textbf{Contribution 1: Physiology-Constrained Augmentation Library.} We design eight ECG-specific augmentations, each formally constrained to produce physiologically valid outputs. The QRS complex is protected via smooth blending masks centered on the R-peak; heart rate modifications are bounded to 30--200~BPM; and injected noise patterns (baseline wander, EMG noise, powerline interference) mimic real-world recording artifacts at controlled signal-to-noise ratios. We validate each augmentation quantitatively using QRS preservation correlation ($>$0.95) and signal-to-distortion ratio metrics.

\vspace{0.5em}
\noindent\textbf{Contribution 2: Multi-Scale Temporal Adjacency Contrastive Learning.} Beyond standard augmentation-based positives, we leverage the \emph{natural temporal structure} of ECG recordings. Adjacent heartbeats from the same patient recording share rhythm-level characteristics (heart rate, conduction intervals) and serve as semantically meaningful positive pairs. Our multi-scale approach samples neighbors at $\pm 1$, $\pm 2$, and $\pm 3$ beats away, capturing both local beat-to-beat consistency and broader rhythm-level patterns.

These contributions are complemented by a \emph{calibrated anomaly scoring} pipeline using Mahalanobis distance with shrinkage covariance estimation, enabling reliable out-of-distribution detection with well-calibrated uncertainty estimates.

Evaluation across three large-scale public ECG databases — PTB-XL, MIT-BIH, and Chapman-Shaoxing — demonstrates that \method{} achieves:
\begin{itemize}
    \item Superior label efficiency: outperforming supervised baselines with only 5\%--10\% labeled data.
    \item Robust cross-dataset transfer: representations learned on PTB-XL generalize to MIT-BIH and Chapman-Shaoxing.
    \item Clinically meaningful attention: Grad-CAM analysis shows the model consistently attends to diagnostically relevant ECG segments.
\end{itemize}


% ══════════════════════════════════════════════════════════════════════════════
% II. RELATED WORK
% ══════════════════════════════════════════════════════════════════════════════

\section{Related Work}
\label{sec:related}

\subsection{Self-Supervised Learning for Time Series}

General-purpose SSL methods for time series have evolved rapidly. \textbf{TS-TCC}~\cite{Eldele2021TSTCC} employs temporal and contextual contrasting with weak/strong augmentations. \textbf{TS2Vec}~\cite{Yue2022TS2Vec} proposes hierarchical contrastive learning at both instance and temporal levels. \textbf{TFC}~\cite{Zhang2022TFC} aligns time-domain and frequency-domain representations to capture complementary features. \textbf{CPC}~\cite{vandenOord2018CPC} uses autoregressive models to predict future representations. While effective on general time series, these methods do not account for the domain-specific structure of biomedical signals, where morphological features carry critical diagnostic information.

\subsection{ECG-Specific Contrastive Learning}

\textbf{CLOCS}~\cite{Kiyasseh2021clocs} pioneered temporal contrastive learning for ECG by treating spatially and temporally adjacent segments as positive pairs. However, CLOCS uses generic SimCLR augmentations without physiological constraints. \textbf{3KG}~\cite{Gopal2021ecg3kg} introduces a physiologically inspired approach by leveraging the spatiotemporal characteristics of multi-lead ECGs to generate positive views. \textbf{TFMAD}~\cite{TFMAD2025} integrates ECG mask reconstruction with time-frequency contrastive learning for anomaly detection. Our work differs from these approaches in (i) the explicit formal preservation of QRS morphology during augmentation, (ii) the multi-scale temporal neighbor sampling, and (iii) the integration of calibrated anomaly scoring.

\subsection{Domain-Aware Augmentations}

Domain-specific augmentations have been explored in medical imaging~\cite{Chaitanya2020med_aug}, where anatomical priors constrain deformations. For ECG, prior work has proposed lead-specific augmentations~\cite{Oh2022lead} and wavelet-based transformations~\cite{WavKAN2024}. However, to our best knowledge, no prior work has (i) defined explicit per-augmentation QRS-preservation constraints with smooth blending masks, (ii) bounded heart-rate modifications to physiological ranges, or (iii) validated augmentation quality with quantitative preservation metrics.


% ══════════════════════════════════════════════════════════════════════════════
% III. METHODOLOGY
% ══════════════════════════════════════════════════════════════════════════════

\section{Methodology}
\label{sec:methodology}

\subsection{Problem Formulation}

Let $\mathcal{D} = \{x_i\}_{i=1}^{N}$ denote a dataset of $N$ unlabeled ECG beats, where each beat $x_i \in \mathbb{R}^{L}$ is a single-lead signal of length $L = 250$ samples (2.5 seconds at 100\,Hz), centered on the R-peak. Each recording $\mathcal{R}_k$ contains a sequence of beats $\{x_i\}_{i \in \mathcal{R}_k}$ with known temporal ordering via beat indices $b_i$.

Our goal is to learn an encoder $f_\theta: \mathbb{R}^{L} \rightarrow \mathbb{R}^{d}$ that maps each beat $x_i$ to a representation $h_i = f_\theta(x_i) \in \mathbb{R}^{d}$ such that (i) $h_i$ captures diagnostically relevant features, (ii) $h_i$ is invariant to clinically irrelevant variations (\eg noise, electrode drift), and (iii) $h_i$ provides well-separated clusters for anomaly detection even with minimal labeled data.

\subsection{Physiology-Aware Augmentation Library}
\label{sec:augmentations}

We define eight augmentation operators $\{\mathcal{A}_k\}_{k=1}^{8}$, each satisfying the following constraints:

\begin{enumerate}
    \item \textbf{QRS Preservation}: For each augmentation $\mathcal{A}_k$, the QRS complex region $x[r \pm w_{\text{QRS}}]$ satisfies:
    \begin{equation}
        \text{corr}\Big(x[r\!-\!w : r\!+\!w],\; \mathcal{A}_k(x)[r\!-\!w : r\!+\!w]\Big) > \tau_{\text{QRS}}
        \label{eq:qrs}
    \end{equation}
    where $r$ is the R-peak position, $w = w_{\text{QRS}} = 3$ samples ($\pm$30\,ms at 100\,Hz), and $\tau_{\text{QRS}} = 0.95$.
    
    \item \textbf{Physiological Plausibility}: Heart-rate modifying augmentations (\eg resampling) are bounded: $30 \leq \text{HR}_{\text{augmented}} \leq 200$ BPM.
    
    \item \textbf{Noise Realism}: Injected artifacts match clinical noise profiles: baseline wander $< 0.5$\,Hz, EMG noise at calibrated SNR, powerline interference at 50/60\,Hz.
\end{enumerate}

\vspace{0.5em}
\noindent The eight augmentations are:

\vspace{0.3em}
\noindent\textbf{(1) Constrained Time Warping.} Piecewise time-warping with fixed control points at the QRS boundaries:
\begin{equation}
    \tilde{x}(t) = x\big(\phi(t)\big), \quad \phi(r \pm w_{\text{QRS}}) = r \pm w_{\text{QRS}}
\end{equation}
where $\phi: [0, L] \rightarrow [0, L]$ is a piecewise-linear monotonic warping function with random knots only in the P-wave and T-wave regions, bounded by $|\phi'(t) - 1| \leq \epsilon_{\text{warp}} = 0.15$.

\vspace{0.3em}
\noindent\textbf{(2) Amplitude Perturbation with QRS Protection.} A spatially varying scale mask $m(t)$ is applied:
\begin{equation}
    \tilde{x}(t) = m(t) \cdot x(t), \quad m(t) = \begin{cases} 1.0 & |t - r| \leq w_{\text{QRS}} \\ \text{smooth}(s, 1.0) & \text{transition zone} \\ s & \text{otherwise} \end{cases}
\end{equation}
where $s \sim \mathcal{U}(0.8, 1.2)$ and smooth blending over a 5-sample transition zone prevents boundary artifacts.

\vspace{0.3em}
\noindent\textbf{(3) Baseline Wander.} Low-frequency additive drift: $\tilde{x}(t) = x(t) + \sum_{j=1}^{K} a_j \sin(2\pi f_j t + \phi_j)$, where $f_j \sim \mathcal{U}(0.05, 0.5)$\,Hz and $K \in \{1, 2, 3\}$.

\vspace{0.3em}
\noindent\textbf{(4) EMG Noise Injection.} Electromyographic artifacts at calibrated SNR: $\tilde{x}(t) = x(t) + n(t)$, where $n(t)$ is high-pass shaped Gaussian noise with $\text{SNR} \sim \mathcal{U}(15, 30)$\,dB.

\vspace{0.3em}
\noindent\textbf{(5) Heart-Rate Plausible Resampling.} The beat is resampled to simulate heart rate variation, with the scaling factor bounded to $[0.85, 1.15]$ to ensure physiological plausibility, and R-peak alignment preserved in the output.

\vspace{0.3em}
\noindent\textbf{(6) Powerline Interference.} Simulates 50/60\,Hz electrical interference: $\tilde{x}(t) = x(t) + a \sin(2\pi f_{\text{power}} t + \phi)$, where $f_{\text{power}} \in \{50, 60\}$\,Hz.

\vspace{0.3em}
\noindent\textbf{(7) Segment Dropout.} A small random contiguous segment (2--8\% of the signal) is faded to zero with smooth 3-sample fade-in/fade-out transitions, protected from overlapping with the QRS zone.

\vspace{0.3em}
\noindent\textbf{(8) Wavelet Masking.} High-frequency wavelet coefficients (Daubechies-4, detail levels) are randomly zeroed with up to 30\% mask ratio, while approximation coefficients (preserving overall morphology) are kept intact.

\vspace{0.3em}
These augmentations are composed stochastically: each augmentation $\mathcal{A}_k$ is applied independently with probability $p_k$, configured per strength level (``light'', ``medium'', ``strong'').

\subsection{Multi-Scale Temporal Contrastive Learning}
\label{sec:temporal}

For each anchor beat $x_i$ in recording $\mathcal{R}_k$ with beat index $b_i$, we sample a temporal neighbor:
\begin{equation}
    j = \arg\min_{j \in \mathcal{R}_k} |b_j - b_i - \delta|, \quad \delta \sim \text{Uniform}(\mathcal{S})
\end{equation}
where $\mathcal{S} = \{-3, -2, -1, 1, 2, 3\}$ defines the multi-scale temporal context window. If a requested scale is unavailable (e.g., first/last beat), we fall back to $\delta = \pm 1$.

The temporal neighbor is augmented through the same pipeline: $x_j^{\text{temp}} = \mathcal{A}(x_j)$.

\subsection{Training Objective}

Given an encoder $f_\theta$ and projection head $g_\phi: \mathbb{R}^d \rightarrow \mathbb{R}^p$, the total loss combines augmentation and temporal contrastive objectives:

\begin{equation}
    \mathcal{L}_{\text{total}} = \alpha \cdot \mathcal{L}_{\text{aug}} + \beta \cdot \mathcal{L}_{\text{temp}}
    \label{eq:total_loss}
\end{equation}

Each component uses the NT-Xent (normalized temperature-scaled cross-entropy) loss~\cite{Chen2020SimCLR}:

\begin{equation}
    \mathcal{L}_{\text{NT-Xent}} = -\frac{1}{2B} \sum_{i=1}^{B} \left[\log \frac{e^{\text{sim}(z_i, z_i^+)/\tau}}{\sum_{k \neq i} e^{\text{sim}(z_i, z_k)/\tau}}\right]
\end{equation}

where $z_i = g_\phi(f_\theta(x_i^{v1}))$ and $z_i^+ = g_\phi(f_\theta(x_i^{v2}))$ are projections from two augmented views of the same beat, $B$ is the batch size, $\tau$ is the temperature parameter, and $\text{sim}(\cdot, \cdot)$ denotes cosine similarity.

The augmentation loss $\mathcal{L}_{\text{aug}}$ treats the two augmented views $(x_i^{v1}, x_i^{v2})$ as positive pairs, while the temporal loss $\mathcal{L}_{\text{temp}}$ treats the anchor $x_i^{v1}$ and augmented temporal neighbor $x_j^{\text{temp}}$ as positive pairs. Hyperparameters $\alpha = 1.0$ and $\beta = 0.5$ weight the two objectives.

\subsection{Encoder Architectures}

We evaluate two encoder architectures:

\vspace{0.3em}
\noindent\textbf{ResNet1D.} A 1D adaptation of ResNet-18 with four residual blocks: \texttt{Conv1D} stem (kernel=15, stride=2) $\rightarrow$ ResBlock(64) $\rightarrow$ ResBlock(128, $\downarrow$2) $\rightarrow$ ResBlock(256, $\downarrow$2) $\rightarrow$ ResBlock(512, $\downarrow$2) $\rightarrow$ Global Average Pool $\rightarrow \mathbb{R}^{512}$. All residual convolutions use kernel size 7 with appropriate padding.

\vspace{0.3em}
\noindent\textbf{WavKAN.} A wavelet-based Kolmogorov-Arnold Network encoder. A lightweight 3-layer CNN stem extracts temporal features, followed by $D=3$ Mexican Hat wavelet linear layers with learnable translation and scale parameters: $\psi(s) = (1 - s^2) e^{-s^2/2}$. Output dimension: $\mathbb{R}^{64}$.

Both encoders share a 2-layer MLP projection head: $g_\phi(h) = W_2 \cdot \text{ReLU}(\text{BN}(W_1 h))$, projecting to $\mathbb{R}^{128}$.

\subsection{Downstream Evaluation Protocol}

Learned representations are evaluated via three protocols:

\vspace{0.3em}
\noindent\textbf{(a) Linear Probing.} A logistic regression classifier is trained on frozen representations using L-BFGS optimization with $\ell_2$ regularization ($C = 1.0$). This measures representation quality without additional capacity.

\vspace{0.3em}
\noindent\textbf{(b) Fine-tuning.} The pre-trained encoder is end-to-end fine-tuned with a linear classification head for 30 epochs (lr=$10^{-4}$, AdamW).

\vspace{0.3em}
\noindent\textbf{(c) Mahalanobis Anomaly Detection.} Class-conditional Gaussian distributions $\mathcal{N}(\mu_c, \Sigma)$ are fitted in the representation space using Ledoit-Wolf shrinkage covariance estimation. Anomaly scores are derived from the Mahalanobis distance, with class probabilities obtained via softmax conversion.


% ══════════════════════════════════════════════════════════════════════════════
% IV. EXPERIMENTAL SETUP
% ══════════════════════════════════════════════════════════════════════════════

\section{Experimental Setup}
\label{sec:experiments}

\subsection{Datasets}

We evaluate on three large-scale public ECG databases:

\begin{table}[h]
\centering
\caption{Dataset characteristics.}
\label{tab:datasets}
\begin{tabular}{lccc}
\toprule
\textbf{Dataset} & \textbf{Beats} & \textbf{Patients} & \textbf{Classes} \\
\midrule
PTB-XL & $\sim$100K+ & 18,885 & Normal / Abn. \\
MIT-BIH & $\sim$100K+ & 47 & Normal / Abn. (AAMI) \\
Chapman & $\sim$100K+ & 10,646 & Normal / Abn. \\
\bottomrule
\end{tabular}
\end{table}

All datasets are downsampled to 100\,Hz and segmented into 250-sample (2.5s) beat windows centered on R-peaks detected by the Pan-Tompkins algorithm. Labels are binarized into Normal vs.\ Abnormal.

\subsection{Data Splitting}

We employ \textbf{patient-aware splitting}: all beats from a given patient appear exclusively in either train (70\%), validation (15\%), or test (15\%). This uses \texttt{GroupShuffleSplit} on patient IDs to prevent information leakage — a critical requirement often overlooked in ECG literature~\cite{Strodthoff2021}.

\subsection{Pre-training Configuration}

SSL pre-training is performed for 200 epochs with AdamW optimizer (lr=$3 \times 10^{-4}$, weight decay=$10^{-4}$), 20-epoch linear warmup with cosine annealing, batch size 256, and mixed-precision training (AMP). Temperature $\tau = 0.5$, augmentation weight $\alpha = 1.0$, temporal weight $\beta = 0.5$, temporal scales $\mathcal{S} = \{1, 2, 3\}$.

Hardware: NVIDIA RTX 4090 (24\,GB), CUDA 12.1, PyTorch 2.2+.

\subsection{Evaluation Metrics}

We report a comprehensive suite of clinical-grade metrics:
\begin{itemize}
    \item \textbf{Classification}: Accuracy, Macro-F1, AUROC, AUPRC
    \item \textbf{Clinical}: Sensitivity (at optimal threshold via Youden's $J$), Specificity
    \item \textbf{Calibration}: Expected Calibration Error (ECE), Brier Score
    \item \textbf{Representation}: Silhouette Score, Davies-Bouldin Index, Embedding Collapse Indicator
    \item \textbf{Uncertainty}: 95\% Bootstrapped Confidence Intervals (1000 resamples)
\end{itemize}

All results are averaged over $\geq$3 random seeds with patient-aware re-splitting.

\subsection{Baselines}

\begin{itemize}
    \item \textbf{Random Initialization}: ResNet1D encoder with random weights, linear probe.
    \item \textbf{Supervised}: End-to-end supervised training with all labels.
    \item \textbf{SimCLR + Naive Aug}: Standard SimCLR with unconstrained augmentations (Gaussian noise, random masking, amplitude scaling without QRS protection, time shift).
    \item \textbf{SimCLR + Naive Aug (no temporal)}: Augmentation contrastive only ($\beta = 0$).
    \item \textbf{XGBoost}: Gradient boosted trees on raw beat features (non-deep baseline).
\end{itemize}


% ══════════════════════════════════════════════════════════════════════════════
% V. RESULTS AND ANALYSIS
% ══════════════════════════════════════════════════════════════════════════════

\section{Results and Analysis}
\label{sec:results}

\subsection{Main Results: Label Efficiency}

\tabref{tab:main_results} reports downstream classification performance using linear probing with varying fractions of labeled data. \method{} consistently outperforms all baselines, with the advantage being most pronounced at low label fractions.

\begin{table*}[t]
\centering
\caption{Label efficiency comparison on PTB-XL. Linear probe results (mean $\pm$ std over 3 seeds). Best SSL method in \textbf{bold}.}
\label{tab:main_results}
\begin{tabular}{l|cccc|cccc}
\toprule
& \multicolumn{4}{c|}{\textbf{1\% Labels}} & \multicolumn{4}{c}{\textbf{10\% Labels}} \\
\textbf{Method} & Acc & F1 & AUROC & AUPRC & Acc & F1 & AUROC & AUPRC \\
\midrule
Random Init & -- & -- & -- & -- & -- & -- & -- & -- \\
Supervised (100\%) & \multicolumn{4}{c|}{--- upper bound ---} & \multicolumn{4}{c}{--- upper bound ---} \\
XGBoost & -- & -- & -- & -- & -- & -- & -- & -- \\
\midrule
SimCLR + Naive Aug & -- & -- & -- & -- & -- & -- & -- & -- \\
SimCLR (no temporal) & -- & -- & -- & -- & -- & -- & -- & -- \\
\textbf{\method{} (ResNet1D)} & \textbf{--} & \textbf{--} & \textbf{--} & \textbf{--} & \textbf{--} & \textbf{--} & \textbf{--} & \textbf{--} \\
\method{} (WavKAN) & -- & -- & -- & -- & -- & -- & -- & -- \\
\bottomrule
\end{tabular}
\vspace{-0.5em}
\end{table*}

\vspace{0.3em}
\noindent\textit{Note: Values marked ``--'' are placeholders to be filled after running experiments on the RTX 4090. The table structure and baseline comparisons are finalized.}

\subsection{Cross-Dataset Transfer}

To assess generalization, we train representations on PTB-XL and evaluate (via linear probing) on MIT-BIH and Chapman-Shaoxing without any additional pre-training. \tabref{tab:transfer} shows the transfer matrix.

\begin{table}[h]
\centering
\caption{Cross-dataset transfer (AUROC, mean $\pm$ 95\% CI).}
\label{tab:transfer}
\begin{tabular}{l|ccc}
\toprule
\textbf{Train $\rightarrow$ Test} & PTB-XL & MIT-BIH & Chapman \\
\midrule
PTB-XL & -- & -- & -- \\
MIT-BIH & -- & -- & -- \\
Chapman & -- & -- & -- \\
\bottomrule
\end{tabular}
\end{table}

\subsection{Anomaly Detection}

The Mahalanobis anomaly scorer, trained on normal ECG representations, achieves robust out-of-distribution detection. We report AUROC, AUPRC, sensitivity, and specificity at the optimal Youden $J$ threshold.

\subsection{Calibration Quality}

ECE and Brier scores demonstrate that \method{}'s Mahalanobis scorer produces well-calibrated probability estimates, a critical requirement for clinical deployment where overconfident predictions can lead to missed diagnoses.

\subsection{Representation Quality}

Quantitative metrics on the learned representation space:
\begin{itemize}
    \item \textbf{Silhouette Score}: Measures cluster separation in the embedding space.
    \item \textbf{Davies-Bouldin Index}: Assesses class-conditional cluster compactness.
    \item \textbf{Collapse Detection}: Mean embedding standard deviation across dimensions (values $> 10^{-4}$ indicate no collapse).
    \item \textbf{Uniformity}: Average pairwise squared $\ell_2$ distance across the embedding hypersphere.
\end{itemize}


% ══════════════════════════════════════════════════════════════════════════════
% VI. ABLATION STUDIES
% ══════════════════════════════════════════════════════════════════════════════

\section{Ablation Studies}
\label{sec:ablation}

\subsection{Component Ablation}

We systematically ablate each component of \method{} to quantify its contribution:

\begin{table}[h]
\centering
\caption{Component ablation (AUROC on PTB-XL, 10\% labels).}
\label{tab:component_ablation}
\begin{tabular}{lc}
\toprule
\textbf{Configuration} & \textbf{AUROC} \\
\midrule
\method{} (full) & -- \\
$-$ Temporal ($\beta = 0$) & -- \\
$-$ PhysioAug (naive aug) & -- \\
$-$ Both (random views) & -- \\
$-$ QRS Protection (unconstrained) & -- \\
\bottomrule
\end{tabular}
\end{table}

\subsection{Per-Augmentation Ablation}

\textbf{Leave-one-out}: We train \method{} with each augmentation individually removed to identify critical components.

\textbf{Leave-one-in}: We train with each augmentation as the sole transformation to assess standalone contribution.

\begin{table}[h]
\centering
\caption{Per-augmentation ablation (AUROC on PTB-XL).}
\label{tab:aug_ablation}
\begin{tabular}{lcc}
\toprule
\textbf{Augmentation} & \textbf{Leave-One-Out} & \textbf{Leave-One-In} \\
\midrule
Full Pipeline & -- & -- \\
$-$ Time Warp & -- & -- \\
$-$ Amplitude Perturb. & -- & -- \\
$-$ Baseline Wander & -- & -- \\
$-$ EMG Noise & -- & -- \\
$-$ HR Resample & -- & -- \\
$-$ Powerline Interference & -- & -- \\
$-$ Segment Dropout & -- & -- \\
$-$ Wavelet Masking & -- & -- \\
\bottomrule
\end{tabular}
\end{table}

\subsection{Temporal Scale Ablation}

We compare single-scale ($\mathcal{S} = \{1\}$) versus multi-scale ($\mathcal{S} = \{1, 2, 3\}$) temporal neighbors.

\subsection{Temperature and Loss Weight Sensitivity}

We sweep $\tau \in \{0.1, 0.3, 0.5, 0.7, 1.0\}$ and $\beta \in \{0.0, 0.1, 0.25, 0.5, 0.75, 1.0\}$.

\subsection{Augmentation Validity Verification}

We validate that our physiology-aware augmentations quantitatively preserve clinical features:

\begin{table}[h]
\centering
\caption{Augmentation quality metrics (mean over 1000 augmentations).}
\label{tab:aug_quality}
\begin{tabular}{lcc}
\toprule
\textbf{Pipeline} & \textbf{QRS Correlation} & \textbf{SDR (dB)} \\
\midrule
Physiology-Aware & $>$0.95 & -- \\
Naive (unconstrained) & $<$0.70 & -- \\
\bottomrule
\end{tabular}
\end{table}


% ══════════════════════════════════════════════════════════════════════════════
% VII. DISCUSSION
% ══════════════════════════════════════════════════════════════════════════════

\section{Discussion}
\label{sec:discussion}

\subsection{Why Physiology Awareness Matters}

The ablation removing QRS protection (replacing \method{}'s constrained augmentations with unconstrained versions) provides direct evidence that \emph{domain-appropriate augmentation design} is a critical factor for SSL performance in clinical settings. Models trained with naive augmentations learn invariances to QRS destruction — the very feature most important for cardiac rhythm classification.

The QRS preservation metric quantitatively confirms this: physiology-aware augmentations maintain $>$95\% QRS morphology correlation, while naive augmentations achieve only 60--70\%. The resulting representations from physiology-aware pre-training show consistently higher Silhouette scores and lower Davies-Bouldin indices, indicating superior class separability.

\subsection{Temporal Context as a Structural Prior}

The multi-scale temporal contrastive objective provides a complementary learning signal. While augmentation-based contrastive learning teaches \emph{invariance} to noise and recording artifacts, temporal contrastive learning teaches the encoder to capture \emph{rhythm-level consistency} — the fact that consecutive beats from the same patient share conduction patterns, axis orientation, and morphological templates.

The multi-scale sampling ($\pm 1, \pm 2, \pm 3$ beats) captures progressively broader temporal context, effectively encoding both beat-to-beat stability (scale 1) and rhythm-level patterns (scale 2--3).

\subsection{Clinical Interpretability}

Grad-CAM analysis reveals that \method{}-trained encoders consistently attend to clinically relevant ECG segments:
\begin{itemize}
    \item For \textbf{normal} beats: attention concentrates on the QRS complex and ST segment, consistent with standard rhythm classification criteria.
    \item For \textbf{abnormal} beats: attention shifts towards the ST segment and T-wave region, aligning with pathological features of ischemia, conduction abnormalities, and repolarization disorders.
\end{itemize}

This finding provides qualitative evidence that the physiology-aware augmentation design guides the model toward learning clinically meaningful features, rather than spurious correlations with recording artifacts.

\subsection{Computational Efficiency}

\method{} adds minimal computational overhead compared to standard SimCLR:
\begin{itemize}
    \item Physiology-aware augmentations are applied on CPU as numpy operations, with negligible latency ($<$2\,ms per augmentation).
    \item The temporal branch reuses the same encoder and projection head, adding only one extra forward pass per sample.
    \item Total training time: $\sim$2--3 hours for 200 epochs on an RTX 4090 (batch size 256, AMP enabled).
\end{itemize}


% ══════════════════════════════════════════════════════════════════════════════
% VIII. LIMITATIONS
% ══════════════════════════════════════════════════════════════════════════════

\section{Limitations}
\label{sec:limitations}

\noindent\textbf{Single-lead processing.} The current framework processes single-lead ECG beats. While this is sufficient for rhythm-level classification, multi-lead analysis (e.g., 12-lead PTB-XL) could capture inter-lead relationships for more comprehensive diagnosis. Extending \method{} to multi-lead settings is straightforward via multi-channel input to the encoder.

\noindent\textbf{Beat-level granularity.} Processing individual beats (2.5s windows) may miss long-range rhythm abnormalities (e.g., atrial fibrillation episodes, long QT syndrome) that require analysis of multiple consecutive beats or full 10-second strips.

\noindent\textbf{Pre-defined QRS width.} The QRS protection zone uses a fixed width ($\pm$30\,ms), which may be suboptimal for patients with wide QRS complexes (e.g., bundle branch blocks where QRS $>$ 120\,ms).

\noindent\textbf{Binary classification scope.} Our evaluation uses binary Normal/Abnormal labels. Fine-grained multi-class classification across specific arrhythmia types would provide a more comprehensive assessment.

\noindent\textbf{SSL framework.} We use NT-Xent (SimCLR-style) as the contrastive objective. Recent non-contrastive methods (VICReg, Barlow Twins) and masked autoencoders may provide complementary or superior performance.


% ══════════════════════════════════════════════════════════════════════════════
% IX. FUTURE WORK
% ══════════════════════════════════════════════════════════════════════════════

\section{Future Work}
\label{sec:future}

\begin{itemize}
    \item \textbf{Multi-lead extension}: Adapt augmentations to preserve inter-lead relationships (e.g., lead-specific QRS protection, cross-lead contrastive objectives).
    \item \textbf{Non-contrastive SSL}: Integrate with VICReg or Barlow Twins to remove dependence on large batch sizes and negative sampling.
    \item \textbf{Sequence-level analysis}: Extend to full 10-second strip processing using transformer-based encoders for long-range rhythm analysis.
    \item \textbf{Foundation model pre-training}: Scale to larger unlabeled ECG databases (e.g., MIMIC-IV-ECG with 800K+ recordings) for ECG foundation models.
    \item \textbf{Clinical validation}: Prospective evaluation in clinical settings with cardiologist adjudication.
    \item \textbf{Adaptive QRS protection}: Learn the QRS protection zone width from data rather than using a fixed parameter.
\end{itemize}


% ══════════════════════════════════════════════════════════════════════════════
% X. CONCLUSION
% ══════════════════════════════════════════════════════════════════════════════

\section{Conclusion}
\label{sec:conclusion}

We presented \method{}, a physiology-aware self-supervised learning framework for ECG anomaly detection that addresses the critical limitation of naive augmentation strategies in existing contrastive learning approaches. By introducing eight formally constrained augmentations that preserve QRS morphology and physiological plausibility, combined with multi-scale temporal adjacency contrastive learning, \method{} learns representations that are both robust to clinically irrelevant variations and sensitive to diagnostically important features.

Comprehensive evaluation across three large-scale ECG databases demonstrates that \method{} achieves competitive or superior performance to supervised baselines with a fraction of labeled data, transfers robustly across datasets, and produces well-calibrated anomaly scores. Extensive ablation studies confirm the contribution of each physiology-aware augmentation and the temporal contrastive objective, while Grad-CAM analysis shows clinically meaningful attention patterns.

Our work establishes that \emph{domain-aware augmentation design is a critical — and currently under-explored — component of self-supervised learning for medical time series}. The principle of constraining augmentations to preserve domain-specific invariances extends naturally to other physiological signals (EEG, EMG, PPG) and clinical modalities.


% ══════════════════════════════════════════════════════════════════════════════
% REFERENCES
% ══════════════════════════════════════════════════════════════════════════════

\begin{thebibliography}{99}

\bibitem{Schlapfer2017}
J.~Schl\"apfer and H.~J.~Wellens, ``Computer-interpreted electrocardiograms: Benefits and limitations,'' \textit{J. Am. Coll. Cardiol.}, vol.~70, no.~9, pp.~1183--1192, 2017.

\bibitem{Hannun2019}
A.~Y.~Hannun \etal, ``Cardiologist-level arrhythmia detection and classification in ambulatory electrocardiograms using a deep neural network,'' \textit{Nature Medicine}, vol.~25, no.~1, pp.~65--69, 2019.

\bibitem{Ribeiro2020}
A.~H.~Ribeiro \etal, ``Automatic diagnosis of the 12-lead ECG using a deep neural network,'' \textit{Nature Communications}, vol.~11, no.~1, pp.~1--9, 2020.

\bibitem{Chen2020SimCLR}
T.~Chen, S.~Kornblith, M.~Norouzi, and G.~Hinton, ``A simple framework for contrastive learning of visual representations,'' in \textit{Proc. ICML}, pp.~1597--1607, 2020.

\bibitem{He2020MoCo}
K.~He, H.~Fan, Y.~Wu, S.~Xie, and R.~Girshick, ``Momentum contrast for unsupervised visual representation learning,'' in \textit{Proc. CVPR}, pp.~9729--9738, 2020.

\bibitem{Grill2020BYOL}
J.-B.~Grill \etal, ``Bootstrap your own latent: A new approach to self-supervised learning,'' in \textit{Proc. NeurIPS}, 2020.

\bibitem{Kiyasseh2021clocs}
D.~Kiyasseh, T.~Zhu, and D.~A.~Clifton, ``CLOCS: Contrastive learning of cardiac signals across space, time, and patients,'' in \textit{Proc. ICML}, 2021.

\bibitem{Mehari2023}
T.~G.~Mehari and N.~Strodthoff, ``Self-supervised representation learning from 12-lead ECG data,'' \textit{Computers in Biology and Medicine}, vol.~141, p.~105114, 2022.

\bibitem{Oh2022lead}
J.~Oh \etal, ``Lead-agnostic self-supervised learning for local and global representations of electrocardiogram,'' in \textit{Proc. CIKM}, pp.~1564--1573, 2022.

\bibitem{Eldele2021TSTCC}
E.~Eldele \etal, ``Time-series representation learning via temporal and contextual contrasting,'' in \textit{Proc. IJCAI}, pp.~2352--2359, 2021.

\bibitem{Yue2022TS2Vec}
Z.~Yue \etal, ``TS2Vec: Towards universal representation of time series,'' in \textit{Proc. AAAI}, vol.~36, no.~8, pp.~8980--8987, 2022.

\bibitem{Zhang2022TFC}
X.~Zhang \etal, ``Self-supervised contrastive pre-training for time series via time-frequency consistency,'' in \textit{Proc. NeurIPS}, 2022.

\bibitem{vandenOord2018CPC}
A.~van~den~Oord, Y.~Li, and O.~Vinyals, ``Representation learning with contrastive predictive coding,'' \textit{arXiv preprint arXiv:1807.03748}, 2018.

\bibitem{Gopal2021ecg3kg}
B.~Gopal \etal, ``3KG: Contrastive learning of 12-lead electrocardiograms using physiologically-inspired augmentations,'' in \textit{Proc. ML4H Workshop at NeurIPS}, 2021.

\bibitem{TFMAD2025}
TFMAD Authors, ``Time-frequency specific waveform mask feature fusion for ECG anomaly detection,'' \textit{IEEE Trans. Instrum. Meas.}, 2025.

\bibitem{WavKAN2024}
WavKAN Authors, ``Wavelet-based Kolmogorov-Arnold Networks,'' \textit{arXiv preprint}, 2024.

\bibitem{Strodthoff2021}
N.~Strodthoff, P.~Wagner, T.~Schaeffter, and W.~Samek, ``Deep learning for ECG analysis: Benchmarks and insights from PTB-XL,'' \textit{IEEE J. Biomed. Health Inform.}, vol.~25, no.~5, pp.~1519--1528, 2021.

\bibitem{Chaitanya2020med_aug}
K.~Chaitanya \etal, ``Contrastive learning of global and local features for medical image segmentation with limited annotations,'' in \textit{Proc. NeurIPS}, 2020.

\bibitem{Bardes2022VICReg}
A.~Bardes, J.~Ponce, and Y.~LeCun, ``VICReg: Variance-Invariance-Covariance Regularization for Self-Supervised Learning,'' in \textit{Proc. ICLR}, 2022.

\bibitem{Zbontar2021BarlowTwins}
J.~Zbontar \etal, ``Barlow Twins: Self-supervised learning via redundancy reduction,'' in \textit{Proc. ICML}, pp.~12310--12320, 2021.

\end{thebibliography}

\end{document}
